\documentclass{article}

\usepackage{amsmath,amssymb}
\usepackage{xurl}
\usepackage[hidelinks]{hyperref}
\setlength{\emergencystretch}{2em}

% Shared commands for notes in docs/paper-gaps/.

\DeclareUnicodeCharacter{2080}{\ifmmode{}_{0}\else\textsubscript{0}\fi}
\DeclareUnicodeCharacter{2081}{\ifmmode{}_{1}\else\textsubscript{1}\fi}
\DeclareUnicodeCharacter{2082}{\ifmmode{}_{2}\else\textsubscript{2}\fi}
\DeclareUnicodeCharacter{2099}{\ifmmode{}_{n}\else\textsubscript{n}\fi}

\newcommand{\Lean}{\textsc{Lean}}
\ExplSyntaxOn
\NewDocumentCommand{\leanid}{m}{
  \ifmmode
    \textsf{\detokenize{#1}}
  \else
    \group_begin:
    \urlstyle{sf}
    \tl_set:Nn \l_tmpa_tl {#1}
    \tl_replace_all:Nnn \l_tmpa_tl {\_} {_}
    \exp_args:NV \path \l_tmpa_tl
    \group_end:
  \fi
}
\ExplSyntaxOff
\providecommand{\tr}{\operatorname{tr}}
\newcommand{\tnleanIssueBase}{https://github.com/LionSR/TNLean/issues/}
\newcommand{\tnleanPrBase}{https://github.com/LionSR/TNLean/pull/}
\newcommand{\tnleanDocBase}{https://sirui-lu.com/TNLean/docs/}
\newcommand{\ghissue}[1]{\href{\tnleanIssueBase#1}{GitHub issue \##1}}
\newcommand{\ghpr}[1]{\href{\tnleanPrBase#1}{PR \##1}}
\newcommand{\doclink}[1]{\href{\tnleanDocBase#1.html}{\path{#1}}}

% Dirac notation and the identity operator, as in blueprint/src/macros/common.tex.
% \providecommand keeps note-local definitions authoritative.
\usepackage{dsfont}
\providecommand{\ket}[1]{|#1\rangle}
\providecommand{\bra}[1]{\langle #1|}
\providecommand{\braket}[2]{\langle #1 | #2 \rangle}
\providecommand{\ketbra}[2]{|#1\rangle\!\langle #2|}
\providecommand{\Id}{\mathds{1}}
\providecommand{\one}{\mathds{1}}
\providecommand{\C}{\mathbb{C}}
\providecommand{\MN}[1]{M_{#1}(\C)}
\providecommand{\MD}{M_D(\C)}

% Verdict marker: \gapnote{<kind>}{<status>}. Kind is the policy
% classification (clarification, local-correction, scope-restriction,
% unfaithful, false-source, open-gap); status is open, wip (elimination
% underway), resolved, or historical. Machine-read by the site index;
% prints a verdict line.
\newcommand{\gapnote}[2]{\par\noindent\textbf{Verdict:} #1 (#2).\par}


% Fallbacks keep this note compilable when a different command.tex is found first.
\providecommand{\leanid}[1]{\textsf{\detokenize{#1}}}
\providecommand{\gapnote}[2]{\par\noindent\textbf{Verdict:} #1 (#2).\par}

\title{The Net Point in the Leading-Eigenvector Lemma of Landau, Vazirani and Vidick}
\author{}
\date{2026-10-28}

\begin{document}
\maketitle
\gapnote{open-gap}{resolved}

\section{The source statement}

In the size-trimming step of Landau, Vazirani and Vidick~\cite{Landau2015Polynomial}
(arXiv:1307.5143v1, \path{algo53.tex}, lines~464--476), $\mathcal N$ is a
finite net over the boundary contractions of bounded bond dimension, whose
cardinality is polynomial in the chain length (lines~454--455).
Lemma \path{l:leading} (lines~493--495) takes for $X$ the element of
$\mathcal N$ closest to the boundary contraction of a low-energy state $w$,
and Claim \path{c:2} (lines~519--520) bounds the size of the resulting set by
the cardinality of $\mathcal N$.

\section{What is formalized}

The leading-vector theorem is stated for any boundary density matrix $X$
whose trace distance to the boundary contraction of $w$ is at most the
tolerance of the convex program. The closest net element is one such $X$, so
the theorem applies to it. The theorem itself makes no assertion about the net
or its cardinality; when it was first formalized, the covering property of the
net and Claim \path{c:2} were therefore an open gap in the size-trimming step.

\section{Resolution}

The size-trimming theorem closes this gap.\footnote{\leanid{TNLean.MPS.AreaLaw.ChainHamiltonian.exists_size_trimmed_family}
    and \leanid{TNLean.MPS.AreaLaw.ChainHamiltonian.exists_size_trimmed_family_polynomial}
    in \doclink{TNLean/MPS/AreaLaw/ViableSet/SizeTrimming}.}
For every actual bond dimension up to the uniform cut bound it chooses a
boundary net, minimizes the left energy at every feasible net point, and takes
the union of the left columns of the solutions, all before the ground vector is
given. Whenever the input family is viable, the net point closest to the
boundary contraction of a low-energy state supplies, through the leading-vector
theorem, a unit vector witnessing viability of the output family. The output
family has cardinality at most $C L^{p}$ with $p=2(qD)^2$, where $C$ and $p$ do
not depend on the chain length. This is Claim \path{c:2}, with the corrected
net exponent of the trace-norm net and the prefactor $(D+1)D$ recorded in the
note on the constant bond dimension. The remaining restrictions of that theorem,
a nonempty chain with at least two states per site and the library's uniform
cut bound, are recorded in the notes on nontrivial chains and on the constant
bond dimension; none of them concerns the net point.

\bibliographystyle{alpha}
\bibliography{references}

\end{document}
